\chapter{代数的演进。}

在数学中，很少有比合成法则更原始的概念：它似乎与整数和可计量的量的计算的最初萌芽不可分割。我们从埃及人和巴比伦人数学留存下来的最古老的文献可以看出，他们已经掌握了一整套关于整数 \textgreater{} 0、有理数 \textgreater{} 0、长度和面积的计算法则；尽管流传到我们手中的巴比伦文本只处理其中数据取明确数值的问题，\(^{1}\) 它们却毫不掩饰地表明，所使用的法则被赋予了一般性，并且在处理一次和二次方程时显示出相当可观的技术娴熟（{[}232{]}，第~179~页及以下）。无论如何，在其中找不到对所用法则的任何论证的迹象，甚至找不到对所出现运算的精确定义：这一切连同那些法则都停留在经验主义的领域之内。

另一方面，类似的关切在古典时代的希腊人那里已经表现得十分清楚；诚然，那时还见不到整数理论的一种公理化处理（这样的公理化直到 19 世纪末才会出现；见第~24~页）；但在欧几里得的《原本》中，已有许多段落对一些其直观程度与整数计算法则相当的“明显”计算法则（例如两个有理数乘积的交换性）给出形式的证明。这类证明中最值得注意的，是与量的理论相关的那些——量的理论是希腊数学最独创的创造（如我们所知，它等价于我们的实数 \textgreater{} 0 的理论；见第~148~页以下）：欧几里得在其中考虑了两个量的比之积等对象，证明了它不依赖于给出这两个比的形式（这是合成法则对等价关系的“商”的第一个例子），并证明了它是交换的（{[}107{]}，第 V 卷，命题 22-23）。\(^{2}\)

不过我们不必讳言，在欧几里得那里，这一朝向严格性的进步，就代数计算的技术而言，是伴随着停滞、甚至某些倒退的。几何学压倒性的优势（量的理论显然正是在几何学的框架内构思出来的）瘫痪了代数记号的一切自主发展：计算中出现的元素必须时时刻刻被“几何地”表示出来；更何况，所出现的两个合成法则不是定义在同一个集合上的（比的加法没有一般性的定义，而两长度之积不是长度而是面积）；由此造成的不灵活，使得处理二次以上的代数关系几乎无法实行。

只是随着古典希腊数学的衰落，人们才看到丢番图回到“计算师”即职业计算者的传统，这些人一直在照原样应用从埃及人和巴比伦人那里继承来的法则：由于不再为他所考虑的那些“数”寻求几何表示，他很自然地被引向发展抽象代数计算的法则；例如，他给出了一些（用现代语言说）等价于公式 \(x^{m+n} = x^{m}x^{n}\)（对 m、n 的小的（正或负）值）的法则（{[}91 a{]}，第 I 卷，第~8-13~页）；稍后一点（第~12-13~页）可以找到“符号法则”的陈述，这是负数计算的最初迹象；最后，丢番图第一次使用一个文字符号来表示方程中的未知数。另一方面，他似乎几乎不去想把他用以解决问题的方法与一般观念联系起来；至于合成法则的公理化概念——它在欧几里得那里已初露端倪——无论对丢番图的思想还是对其直接后继者的思想来说都是陌生的；在代数中，要到 19 世纪初才会重新找到它。

首先需要的是，在其间的几个世纪中，一方面发展出一套适合表达抽象法则的代数记号体系，另一方面，“数”的概念得到拓宽，使人们能够通过对足够多样化的特殊情形的观察上升到一般概念。在这一点上，希腊人创造的量的比的公理化理论是不够的，因为它只是把实数 \textgreater{} 0 的直观概念以及对这些数的运算精确化了，而巴比伦人早已以更含混的形式知道了它们；现在要谈的相反是希腊人未曾设想过的那些“数”，而且起初没有任何可感的“表示”自行呈现出来：一方面是零和负数，它们早在中世纪前期就已出现于印度数学；另一方面是虚数，它们是 16 世纪意大利代数学家的创造。

如果把零撇开——它最初是作为记号中的一个符号被引入的，后来才被当作一个数（见第~45~页）——这些扩张的共同特征是（至少起初）纯粹“形式的”。这应当理解为：新的“数”最初是作为施行于某些情形的运算结果而出现的，在这些情形下，按其严格的定义，它们没有任何意义（例如当 a \textless{} b 时两个整数之差 a - b）：因此人们赋予了它们“假的”、“虚构的”、“荒谬的”、“不可能的”、“想像的”数等等名称。对于一味追求清晰思维的古典时代的希腊人来说，这样的扩张是不可设想的；它们只能出自这样一些计算者：他们比希腊人更倾向于对其方法的力量怀有甚至是几分神秘的信心（用 18 世纪的说法，“分析的普遍性”），任凭自己被计算的机制拖着走而不逐步检验其可靠性；而这种信心事后大多得到了辩护，因为把原本仅对先前已知的数严格有效的计算法则推广到这些新数学对象上去，导出了正确的结果。这就解释了，人们何以越发大胆地就其本身（独立于对具体计算的一切应用）来考虑数概念的这些推广——起初它们只是运算序列中的中介，其出发点和目标都是真正的“数”；一旦迈出这一步，人们便开始为这些新实体寻找或多或少具体的解释，它们由此获得了在数学中被引用的权利。\(^{4}\)

在这一点上，印度人已经意识到负数在某些情形下应有的解释（例如商业问题中的一笔债务）。在其后的几个世纪中，随着希腊和印度数学的方法与成果（经由阿拉伯人）传向西方，人们对这些数的使用越来越熟练，并开始有了对它们的另外一些几何的或动力学的“表示”。无论如何，这与代数记号的逐步改进一起，构成了中世纪末期代数中唯一值得注意的进步。

16 世纪初，代数开始了新的腾飞，这要归功于意大利学派的数学家发现了三次方程、随后四次方程的“根式解”（见第~72~页以下）；正是在这个场合，他们尽管心怀厌恶，却可以说是被迫把虚数引入了计算；不过，对这类“不可能的”数的计算的信心逐渐建立起来，如同对负数一样，尽管在两个多世纪里没有人想出对它们的一种“表示”。

另一方面，代数记号从韦达和笛卡尔那里获得了决定性的改进；从后者开始，代数记号除去少数例外，已经就是我们今天所使用的记号了。

从 17 世纪中叶到 18 世纪末，无穷小演算的创立所敞开的广阔视野，看来在某种程度上造成了对代数总的忽视，特别是对合成法则、对实数和复数本性的数学思考的忽视。\(^{5}\) 于是，力学中早在 17 世纪末就已熟知的力的合成与速度的合成，对代数没有产生任何影响，尽管它们已经以胚胎的形式包含着向量演算。的确，我们要一直等到 1800 年前后导致复数的几何表示的那场思想运动（见第~161~页），才在纯粹数学中看到向量加法的使用。\(^{6}\)

正是在这个时候，合成法则的概念第一次在代数中沿两个不同的方向得到扩展，扩展到这样一些元素：它们与“数”（这个词至此被赋予的最广意义）之间只剩遥远的类比。第一次扩展应归功于 C. F. 高斯，出自他关于整系数二次型 \(ax^{2} + bxy + cy^{2}\) 的数论研究。拉格朗日曾在同一判别式的型的集合中定义了一个等价关系，\(^{7}\) 并且另外证明了一个恒等式，它在该集合中给出一个交换的（并非处处有定义的）合成法则；从这些结果出发，高斯证明这个法则与一个稍有不同的等价关系相容（{[}124 a{]}，第 I 卷，第~272~页）：“由此可以看出”，他当时说，“由两个或多个类合成的类应当理解为什么”。随后他进而研究他刚刚定义的“商”法则，实质上确立了它（用现代语言说）是一条阿贝尔群法则，而且他所用的论证在多数情形下远比高斯所考虑的特殊情形更具普遍性（例如，他用来证明对称元唯一性的论证可以应用于任何合成法则（同上，第~273~页））。但他没有止步于此：稍后重回这一问题时，他认识到类的合成与整数模一个素数的乘法之间的类似\(^{8}\)（同上，第~371~页），但也注意到给定判别式的二次型的类群并不总是循环群；他在这方面给出的指示使我们想到，他或许至少在这一特殊情形下已经认识到了有限阿贝尔群的一般结构（{[}124 a{]}，第 I 卷，第~374~页和第 II 卷，第~266~页）。

我们要谈的另一系列研究同样以群的概念告终，从而以决定性的方式把它引入数学：这就是“置换理论”，它是拉格朗日、范德蒙德和高斯关于代数方程求解的思想的发展。我们不打算在这里给出这一问题的详细历史（见第~75~页以下）；我们应当记住鲁菲尼 {[}265{]}、随后是柯西（{[}56 a{]}，(2)，第 I 卷，第~64~页）对有限集的两个置换之“积”的定义，\(^{9}\) 以及关于有限置换群的最初概念：传递性、本原性、中性元、交换的元素，等等。但这些最初的研究整体上仍相当表面，而真正应当被视为这一理论开创者的是埃瓦里斯特·伽罗瓦：在他那部令人难忘的著作 {[}123{]} 中，他把代数方程的研究化归为他与这些方程相联系的置换群的研究，并对后者作了重大的深化，既涉及群的一般性质（正是伽罗瓦第一个定义了正规子群的概念并认识到它的重要性），也涉及具有特殊性质的群的确定（他在这一方面所得的结果至今仍跻身于该理论最精深的结果之列）。同样也是

伽罗瓦最先有了“群的线性表示”这一想法，\(^{10}\) 而这一事实清楚地表明，他已掌握两个群结构同构的概念，而与其“实现”无关。

然而，即使高斯和伽罗瓦的杰出方法无疑已把他们引向了十分宽广的合成法则概念，他们却未曾有机会专门展开他们在这一点上的想法，其工作对抽象代数的演进也没有产生任何直接的作用。\(^{11}\) 朝向抽象的最明显的进展是在第三个方向上取得的：继关于虚数本性的思考（其几何表示在 19 世纪初曾引起相当可观的工作）之后，英国学派的代数学家从 1830 年到 1850 年首先提炼出合成法则的抽象概念，并立即把这概念应用于一大堆新的数学对象，从而扩大了代数的领域：布尔那里的逻辑代数（见第~8~页）、哈密顿那里的向量、四元数和一般的超复系统 {[}145 a{]}、凯莱那里的矩阵和非结合法则（{[}58{]}，第 I 卷，第~127~页和第~301~页，及第 II 卷，第~185~页和第~475~页）。一场平行的演变在欧洲大陆独立发生，尤其在向量计算（默比乌斯、贝拉维蒂斯）、线性代数和超复系统（格拉斯曼）方面（见第~61~页）。\(^{12}\)

从使代数在 19 世纪上半叶重新焕发活力的这一切独创而多产的思想的沸腾之中，代数脱颖而出，连其趋向都为之一新。在此之前，方法和结果都围绕着代数方程（或数论中的丢番图方程）求解这一中心问题运转：“代数”，塞尔在其《高等代数教程》的引言中说 {[}284{]}，“严格说来就是方程的分析”。1850 年以后，尽管代数教材在很长时间里仍把首要地位留给方程理论，新的研究已不再为求数值方程之解的立竿见影的应用所支配，而是越来越转向我们今天视为代数本质问题的东西，即对代数结构本身的研究。

这项工作相当清晰地分成三股潮流，它们分别延伸着我们上面分析过的三场思想运动，并且彼此平行地推进、相互间没有可察觉的影响，直至 19 世纪的最后几年。\(^{13}\)

首先是德国学派（狄利克雷、库默尔、克罗内克、戴德金、希尔伯特）对代数数理论的建构，它脱胎于高斯的工作——最早的一类研究，即对数 \(a + bi\)（\(a\)、\(b\) 为有理数）的研究，正应归功于高斯。我们不打算在此追述这一理论的演变：我们只需为了自己的目的拣出在其中浮现的抽象代数概念。早在高斯的第一批后继者那里，（代数数的）域的概念就已处于上述全部工作的基础地位（阿贝尔和伽罗瓦关于代数方程的研究也是如此）；当戴德金和韦伯 {[}80{]} 把单变量代数函数论移植到代数数论上时，其应用范围进一步扩大。理想概念的引入也应归功于戴德金（{[}79{]}，第 III 卷，第~251~页），它给出了元素集合之间合成法则的一个新例子；阿贝尔群和模在代数域理论中所起的越来越大的作用，应追溯到他和克罗内克；我们将在后文回到这一点（见第~63~页以下和第~93~页以下）。

我们把线性代数（第~57~页）和超复系统（第~117~页）发展史的介绍也留到以后；它们在 19 世纪末和 20 世纪初的展开没有引入任何新的代数概念：在英国（西尔维斯特、W. 克利福德）和美国（B. 与 C. S. 皮尔斯、迪克森、韦德伯恩），是沿着哈密顿和凯莱开辟的道路；在德国（魏尔斯特拉斯、戴德金、弗罗贝尼乌斯、莫利恩）和法国（拉盖尔、E. 嘉当），则以独立于盎格鲁-撒克逊人的方式并采用颇为不同的方法。

至于群论，它起初主要以有限置换群理论的形式得到发展，这是在伽罗瓦著作发表、并经塞尔的著作 {[}284{]} 传播之后，尤其要归功于 C. 若尔当的巨著《置换论》{[}174 a{]}。若尔当在此书中对前人关于置换群特殊性质（传递性、本原性等）的工作作了极大的完善和总结，所得结果中的大部分至今几乎未被超越；他还深入研究了非常重要的特殊群，即线性群及其子群（见第~138~页）；同样，正是他引入了一个群由另一个群来表示这一基本概念，以及（稍后一点）商群的概念，并且证明了以“若尔当-赫尔德定理”之名著称的定理的一部分。\(^{14}\) 最后，无穷群的最早研究也应归功于若尔当 {[}174 b{]}，几年之后，S. 李与 F. 克莱因和 H. 庞加莱分别沿两个不同的方向对它作了重大的发展。

在此期间，人们逐渐看清：群的本质在于其合成法则，而不在于组成群的那些对象的性质（例如见 {[}58{]}，第 II 卷，第~123~页和第~131~页，以及 {[}79{]}，第 III 卷，第~439~页）。然而，甚至关于有限抽象群的研究，在很长时间里仍被当作关于置换群的研究来构思，直到 1880 年前后，才有了有限群自主理论的自觉发展。我们无法进一步追述这一理论的历史；我们只限于提及至今仍在有限群研究中位居最常用工具之列、而且都可上溯到 19 世纪的两种工具：其一是西罗 \(^{15}\) 关于 p 群的诸定理，其年代为 1872 年 {[}303{]}；其二是由弗罗贝尼乌斯在本世纪最后几年创立的特征标理论（{[}119{]}，第 III 卷，第~1-37~页）。愿意深入了解有限群理论及其所引出的众多困难问题的读者，可参阅专著 {[}44 b{]}、{[}131{]}、{[}291{]} 和 {[}341{]}。

同样是在 1880 年前后，群的表现开始得到系统的研究；在此之前，人们研究的只是特殊群的表现，例如交错群 \(\mathfrak{A}_5\) 在哈密顿著作 {[}145 b{]} 中的表现，或微分方程与黎曼曲面的单值群的表现（施瓦茨、克莱因、施莱夫利）。W. 迪克第一个 {[}97{]} 定义了（但未命名）由有限多个生成元生成的自由群，但他对它本身的兴趣，不如把它当作一个“普适”对象的兴趣——该对象使“由生成元和关系给出的群”这一说法获得确切的定义。迪克发展了凯莱最先提出的一个想法，而且显然还受到上文提到的黎曼学派工作的影响，他描述了具有给定表现的群的一种解释，其中每个生成元由关于两个相切圆或相交圆的反演之积来表示（另见 {[}44 b{]}，第 XVIII 章）。稍后，在庞加莱及其后继者发展了自守函数理论、庞加莱又引入了代数拓扑的工具之后，基本群的首批研究将与群表现的研究并肩推进（德恩、蒂策），两种理论相互支撑。而且恰恰是一位拓扑学家 J. 尼尔森，于 1924 年在其性质的第一个深入研究 {[}234{]} 中引入了“自由群”这一名称；紧接着，E. 阿廷（仍然与拓扑问题相关）引入了群的自由积概念，O. 施赖尔则更一般地定义了带融合子群的自由积。这里我们同样无法处理这一方向后续发展的历史，欲知详情，请参看著作 {[}216{]}。

这里也不是谈论群的概念（以及与之紧密相连的不变量概念）自 19 世纪末以来在分析、几何、力学和理论物理中所享有的异常繁荣的地方。正是以类似的方式，这一概念以及与之相关的一些代数概念（带算子群、环、理想、模）侵入了代数中直到那时看来还相当远离其本域的那些部分——我们在此追述的演化的最后一个时期即以此著称，而且它以我们上文考察过的三种倾向的合流而告终。这一统一首先是现代德国学派的功绩：代数的公理化工作发端于 19 世纪末的戴德金和希尔伯特，E. 施泰尼茨 {[}294 a{]} 强有力地推进了它，随后从 1920 年起，又在 E. 阿廷、E. 诺特及其学派的代数学家（哈塞、克鲁尔、O. 施赖尔、范德瓦尔登）的推动下继续进行。范德瓦尔登出版于 1930 年的教程 {[}317 a{]} 第一次把这些成果汇集在一个完整的阐述之中，为抽象代数此后多种多样的研究开辟了道路并充当了向导。
% semantic-fragments: legacy-input-seam
\relax{}
